% !PRINTED 187
\weberpage{187}
%<<CONT
of an arbitrary point fixed in space, referred to this
system of axes in the unit position $E$, one obtains the coordinates of the same point, referred to the system of axes in the
position $A$ or $B$, by two orthogonal substitutions $A(x)$
and $B(x)$.

Accordingly $AB(x)$ are the coordinates, referred to the
system $A$, of a point $y$ with the coordinates $(y) = B(x)$
in the system $E$. The point $y$ thus has in the system $E$ the same
coordinates as the point $x$ in the system $B$; that is, $y$ lies with respect to the
unit position as $x$ does to the position $B$. If we now carry out the
rotation that leads from $E$ to $B$ in such a way that we hold the point $x$
fast, but let the point $y$ and the system of axes $A$ accompany
the rotation, then the point $y$ arrives at $x$, and the coordinates of $x$, referred to the new situation of the system $A$, are
the same as those of the point $y$ with respect to the original
situation of $A$, that is, $AB(x)$. But this new situation of the system $A$
can evidently also be reached by performing the rotation
that leads to the position $A$ not from the unit position,
but from the position $B$.

From the group of the proper ones one obtains the totality
of all orthogonal substitutions by composition with
\gesperrt{one} improper one, say with $(x_1, x_2, x_3) = (y_1, y_2, - y_3)$,
which, according to the second geometric conception, is a reflection
in the plane $x_1, x_2$, in which the point $y$ is the mirror image
of the point $x$.

Of special interest now are the \gesperrt{finite groups} contained in the totality
of the proper orthogonal substitutions, into which we shall enter more closely later.
Here we wish only to remark the following about the geometric side of this
question. To such a finite group $G$
of degree $g$ of orthogonal substitutions there corresponds a
finite group of positions of a body. If one imagines
the body fixed in these various positions simultaneously
and the whole united into a new rigid body, one
obtains a figure which can cover itself in a finite number $g$ of
distinct positions. For every
group one can find bodies of infinitely many different shapes.
The most intuitive and best-known relations result if one assumes the body bounded by plane faces.
